\documentclass[../DoAn.tex]{subfiles}
\begin{document}
Phụ lục này đặc tả các use case không được trình bày trong mục~\ref{section:2.3}. Mã quy tắc BR tham chiếu tới Bảng~\ref{tab:br}.

\section{Đặc tả use case ``Hủy đặt phòng''}
\begin{dacta}{Đặc tả use case UC06 -- Hủy đặt phòng}{tab:uc06}
Tác nhân & User (chủ booking) \\ \hline
Tiền điều kiện & Booking ở trạng thái PENDING hoặc APPROVED và chưa đến giờ bắt đầu. \\ \hline
Luồng chính &
\begin{buoc}
\item User mở ``Đặt phòng của tôi'', chọn booking và bấm hủy (rút yêu cầu).
\item Hệ thống kiểm tra người gọi là chủ booking và booking còn hoạt động.
\item Với booking APPROVED, hệ thống kiểm tra còn ít nhất 30 phút (BR-08).
\item Hệ thống chuyển CANCELLED, ghi thời điểm thu hồi cho mã PIN nếu có.
\item Hệ thống thông báo ``Yêu cầu đã được hủy'' cho Admin.
\end{buoc} \\ \hline
Luồng phát sinh & 3a. Quá hạn hủy: báo ``Yêu cầu đã duyệt chỉ được hủy trước ít nhất 30 phút''. 2a. Phiên đã bắt đầu: báo không thể hủy. \\ \hline
Hậu điều kiện & Booking CANCELLED, phòng trở lại trống trong khung giờ đó, mã PIN chuyển REVOKED. \\
\end{dacta}

\section{Đặc tả use case ``Hẹn lịch nhận khóa và ủy quyền nhận hộ''}
\begin{dacta}{Đặc tả use case UC07 -- Hẹn lịch nhận khóa và ủy quyền nhận hộ}{tab:uc07}
Tác nhân & Admin, User \\ \hline
Tiền điều kiện & Booking APPROVED tại phòng PHYSICAL\_KEY, chưa đến giờ sử dụng. \\ \hline
Luồng chính &
\begin{buoc}
\item Admin chọn booking, nhập ngày, giờ và địa điểm nhận khóa/thẻ (ví dụ ``Phòng trực A1'').
\item Hệ thống kiểm tra lịch hẹn ở tương lai và trước giờ sử dụng, địa điểm không rỗng (BR-12); lưu và thông báo ``Đã có lịch nhận khóa'' cho User.
\item Nếu không tự đến nhận, User nhập họ tên và mã sinh viên/cán bộ của người nhận hộ.
\item Hệ thống lưu ủy quyền và thông báo cho Admin; người nhận hộ xuất trình thẻ khi nhận.
\end{buoc} \\ \hline
Luồng phát sinh & 2a. Phòng dùng khóa số: báo ``Phòng này không dùng khóa cơ/thẻ''. 3a. Thiếu họ tên hoặc mã: báo lỗi. \\ \hline
Hậu điều kiện & Booking có \texttt{keyPickupAppointment} và tùy chọn \texttt{pickupDelegate}. \\
\end{dacta}

\section{Đặc tả use case ``Quản lý lịch sử dụng và bảo trì''}
\begin{dacta}{Đặc tả use case UC08 -- Quản lý lịch sử dụng và bảo trì}{tab:uc08}
Tác nhân & Admin \\ \hline
Tiền điều kiện & Admin đã đăng nhập. \\ \hline
Luồng chính &
\begin{buoc}
\item Admin mở ``Lịch phòng và bảo trì'', xem lịch các booking đã duyệt.
\item Admin chọn tầng, chọn phòng, nhập ngày, giờ bắt đầu, giờ kết thúc và lý do.
\item Hệ thống kiểm tra phòng tồn tại, khoảng giờ hợp lệ, lý do không rỗng, không trùng bảo trì khác và không trùng booking đã duyệt.
\item Hệ thống lưu lịch bảo trì; phòng hiển thị màu bảo trì trên sơ đồ trong khung giờ đó.
\end{buoc} \\ \hline
Luồng phát sinh & 3a. Trùng booking đã duyệt: báo ``cần đổi phòng trước khi lên lịch bảo trì''. 4a. Admin hủy lịch bảo trì: ghi \texttt{cancelledAt}, phòng khả dụng trở lại. \\ \hline
Hậu điều kiện & Các yêu cầu mới trùng khung giờ bị chặn theo BR-04. \\
\end{dacta}

\section{Đặc tả use case ``Cấu hình nghiệp vụ''}
\begin{dacta}{Đặc tả use case UC09 -- Cấu hình nghiệp vụ}{tab:uc09}
Tác nhân & Admin \\ \hline
Tiền điều kiện & Admin đã đăng nhập. \\ \hline
Luồng chính &
\begin{buoc}
\item Admin mở ``Quy định đặt phòng''; hệ thống hiển thị giá trị hiện tại.
\item Admin sửa số ngày đặt trước tối thiểu, tối đa, số booking hoạt động tối đa mỗi User, hạn hủy, hạn đổi phòng, khoảng đệm mã số, bật/tắt thông báo.
\item Hệ thống kiểm tra các giá trị là số nguyên không âm, tối thiểu không lớn hơn tối đa, giới hạn booking ít nhất bằng 1.
\item Hệ thống lưu cấu hình; các module đọc giá trị mới ở lần kiểm tra tiếp theo.
\end{buoc} \\ \hline
Luồng phát sinh & 3a. Giá trị không hợp lệ: hiển thị thông điệp tương ứng, không lưu. \\ \hline
Hậu điều kiện & Quy tắc BR-03, BR-06, BR-08, BR-09, BR-10 dùng tham số mới. \\
\end{dacta}

\section{Đặc tả use case ``Quản lý SmartLock''}
\begin{dacta}{Đặc tả use case UC10 -- Quản lý SmartLock}{tab:uc10}
Tác nhân & Admin; SmartLock Gateway \\ \hline
Tiền điều kiện & Admin đã đăng nhập; gateway chạy trong cùng mạng. \\ \hline
Luồng chính &
\begin{buoc}
\item Admin mở ``Quản lý khóa'', nhập địa chỉ gateway dạng \texttt{http://<IP>:8135}.
\item Hệ thống chuẩn hóa địa chỉ và gọi \texttt{GET /api/lock/status} để hiển thị trạng thái kết nối OneIoT, lệnh và phản hồi gần nhất.
\item Admin chọn một phòng khóa số để gắn khóa DLWA12.
\item Hệ thống kiểm tra phòng tồn tại và dùng PIN\_CODE, lưu phòng được gắn, người và thời điểm gắn.
\end{buoc} \\ \hline
Luồng phát sinh & 1a. Địa chỉ không bắt đầu bằng http/https: báo lỗi. 3a. Chọn phòng khóa cơ: báo ``Chỉ được gắn SmartLock vào phòng sử dụng khóa mã số''. 4a. Admin tháo khóa: xóa liên kết, mọi yêu cầu tạo mã bị chặn. \\ \hline
Hậu điều kiện & Chỉ booking của phòng đang gắn khóa mới được gửi lệnh tạo mật khẩu. \\
\end{dacta}

\section{Đặc tả use case ``Quản lý tài khoản và quyền theo phòng''}
\begin{dacta}{Đặc tả use case UC11 -- Quản lý tài khoản và quyền theo phòng}{tab:uc11}
Tác nhân & Admin \\ \hline
Tiền điều kiện & Admin đã đăng nhập. \\ \hline
Luồng chính &
\begin{buoc}
\item Admin nhập tên đăng nhập, mật khẩu ban đầu, mã khôi phục và bấm ``Thêm tài khoản''.
\item Hệ thống kiểm tra tên 3--30 ký tự thường, số, dấu chấm, gạch ngang; mật khẩu và mã khôi phục ít nhất 4 ký tự; tên chưa tồn tại.
\item Admin tìm tài khoản, khóa/mở, đặt lại mật khẩu hoặc xóa.
\item Admin xem các phòng mà User được tự tạo mã và thu hồi quyền theo từng phòng.
\end{buoc} \\ \hline
Luồng phát sinh & 3a. Tự khóa hoặc tự xóa tài khoản đang đăng nhập: bị chặn. 3b. Tài khoản còn booking hoạt động: không cho xóa. \\ \hline
Hậu điều kiện & Tài khoản bị khóa không đăng nhập được; quyền bị thu hồi chặn tạo mã theo BR-11. \\
\end{dacta}
\end{document}
